\documentclass[aspectratio=169]{beamer}

\usepackage[spanish,es-tabla]{babel}
\usepackage[utf8]{inputenc}
\usepackage[T1]{fontenc}
\usepackage{booktabs}

\usetheme{Madrid}
\usecolortheme{seahorse}
\setbeamertemplate{navigation symbols}{}
\setbeamertemplate{footline}{%
  \hfill\scriptsize Diapositiva \insertframenumber\hspace{1em}\vspace{0.6em}}

\title{Sistema de visi\'on artificial para conducci\'on aut\'onoma}
\author{Juan Jos\'e Prado Neira \and Javier Rego Bargos}
\date{Exposici\'on final: diapositivas 9--17}

\newcommand{\timing}[2]{%
  \begin{block}{Tiempo de exposici\'on}
  \small #1 segundos \quad $\approx$ \quad #2 palabras a 140 palabras/minuto.
  \end{block}
}

\begin{document}
\setcounter{framenumber}{8}

\begin{frame}{Exportaci\'on e inferencia}
\timing{35}{82}
\begin{itemize}
  \item El modelo entrenado se exporta a ONNX.
  \item Exportaci\'on realizada desde el notebook de entrenamiento.
  \item En runtime se carga desde \texttt{socket\_vision/models/line\_segmentation}.
  \item El sistema ejecuta el modelo ONNX con OpenCV DNN.
  \item Cada frame se normaliza, pasa por la red y se convierte en m\'ascaras binarias.
\end{itemize}
\end{frame}

\begin{frame}{Procesado de m\'ascaras}
\timing{35}{82}
\begin{itemize}
  \item Archivo: \texttt{socket\_vision/vision/perception/line\_segmentation.py}.
  \item Se aplica sigmoide y umbral para obtener las clases.
  \item La m\'ascara del circuito se limpia con operaciones morfol\'ogicas.
  \item Se filtran componentes peque\~nos para reducir ruido.
  \item Se calcula el centro de la zona transitable en una regi\'on de inter\'es.
\end{itemize}
\end{frame}

\begin{frame}{L\'ogica de seguimiento del circuito}
\timing{30}{70}
\begin{itemize}
  \item Se compara el centro del circuito con el centro de la imagen.
  \item Si el circuito aparece desplazado, se genera un error lateral.
  \item Ese error se transforma en \textit{steering}.
  \item Si el coche est\'a centrado, mantiene avance recto.
  \item Si se pierde el circuito, puede pararse por seguridad.
\end{itemize}
\end{frame}

\begin{frame}{Esquiva de cajas}
\timing{40}{93}
\begin{itemize}
  \item Archivo: \texttt{socket\_vision/vision/runtime/autonomy.py}.
  \item Se analiza la m\'ascara de cajas en la zona frontal.
  \item Se calcula posici\'on, \'area y urgencia del obst\'aculo.
  \item Caja a la izquierda: giro hacia la derecha.
  \item Caja a la derecha: giro hacia la izquierda.
  \item Si el riesgo es alto, se reduce velocidad o se detiene.
\end{itemize}
\end{frame}

\begin{frame}{Control del coche en Jetson}
\timing{40}{93}
\begin{itemize}
  \item Archivo: \texttt{socket\_vision/jetson/control\_server.py}.
  \item Recibe comandos \texttt{auto\_drive} por socket.
  \item Aplica \textit{steering} y \textit{throttle} con \texttt{NvidiaRacecar}.
  \item Incluye calibraci\'on de direcci\'on y l\'imites de seguridad.
  \item Permite control manual con prioridad temporal sobre el autom\'atico.
  \item Tiene watchdog para detener el coche si se pierden comandos.
\end{itemize}
\end{frame}

\begin{frame}{Arquitectura general}
\timing{35}{82}
\begin{itemize}
  \item \texttt{circuito\_segmentation}: dataset, entrenamiento, pruebas y exportaci\'on.
  \item \texttt{socket\_vision}: visi\'on en tiempo real, control, comunicaci\'on y Jetson.
  \item El PC se usa para entrenamiento y preparaci\'on.
  \item La Jetson ejecuta captura, comunicaci\'on y control del robot.
  \item El script \texttt{main.sh} sincroniza y lanza el flujo completo.
\end{itemize}
\end{frame}

\begin{frame}{Resultados de prueba real}
\timing{30}{70}
\begin{itemize}
  \item Se valid\'o la integraci\'on completa de visi\'on, comunicaci\'on y control.
  \item El sistema segmenta circuito y cajas en la prueba real.
  \item Calcula el centro de avance y genera comandos de direcci\'on y velocidad.
  \item L\'imite observado: campo de visi\'on reducido o c\'amara demasiado adelantada.
  \item L\'imite observado: latencia por enviar v\'ideo al PC y devolver comandos a Jetson.
\end{itemize}
\end{frame}

\begin{frame}{Conclusiones}
\timing{40}{93}
\begin{itemize}
  \item La segmentaci\'on sem\'antica es adecuada porque da informaci\'on por p\'ixel.
  \item Dos clases son suficientes para la tarea: circuito y caja.
  \item La U-Net ligera permite trabajar con restricciones de c\'omputo.
  \item La l\'ogica de control convierte m\'ascaras en comandos simples y efectivos.
  \item Mejoras clave: c\'amara/posici\'on con m\'as campo de visi\'on y reducci\'on de latencia.
  \item Futuro: m\'as datos, m\'as luz, ajuste de umbrales, DeepLabv3+ o TensorRT.
\end{itemize}
\end{frame}

\begin{frame}{Cierre}
\timing{15}{35}
\begin{itemize}
  \item Se construy\'o una tuber\'ia completa desde captura y etiquetado hasta conducci\'on aut\'onoma.
  \item El resultado no es solo un modelo entrenado: es una integraci\'on completa con Jetson y control del coche.
\end{itemize}
\end{frame}

\end{document}
